Definieren Sie Haftpflichtversicherung und klassifizieren
Sie grob in ihrem rechtlichen Kontext.

Die Haftpflichtversicherung gehört zu den Schadensversicherungen. 
Sie soll den Versicherungsnehmer vor Vermögenseinbußen schützen, die dadurch entstehen, 
dass ihn ein Dritter — zu Recht oder zu Unrecht — auf Schadensersatz in Anspruch nimmt (sog. Passivversicherung). 

Es handelt sich also um eigene Vermögensvorsorge. 
Der Grad zwischen schuldloser und leicht fahrlässiger Verletzung fremder Rechtsgüter ist schmal.
Basis sind (§100-124 VVG) und insbesondere 
§823 BGB (Schadensersatzpflicht der Schaden wird dort als Delikt (anders als Strafbestand) verstanden).

Aus dieser Systematik ergibt sich ein Dreiecksverhältnis, 
wobei zwischen dem Dritten und dem Haftpflichtversicherer grundsätzlich keine unmittelbare Rechtsbeziehung besteht. 

INSERT IMAGE PAGE 83

Häufig ist sie nicht nur eine Eigenschutzversicherung des Versicherungsnehmers, sondern zugleich eine Versicherung für fremde Rechnung dar 
(Ziff. 27 AHB) (z.B. vertretungsberechtigte Personen (Kinder) und Angestellte  bei der Betriebshaftpflichtversicherung (§ 102 I VVG)).

Sie können mit einer Forderungsausfallversicherung kombiniert sein. Das heißt gibt es eine Zahlungs-oder Leistungsunfähigkeit
des schadensersatzpflichtigen Dritten, dann ist der VR leistungspflichtig in dem Umfang wie als hätte der schadensersatzpflichtigen Dritten
eine haftpflichtversicherung beim VR abgeschlossen. Dafür müssen die Ansprüche an den VR abgetreten werden und ein vollstreckbare Ausfertigung
des Urteils ausgehändigt.



Erläuterns Sie das Trennungsprinzip und Bindungswirkung

Elementare Grundlage der Haftpflichtversicherung — und bei der Falllösung stets zu beachten — ist das Trennungsprinzip. 
Es beschreibt die Unterscheidung von Haftung und Deckung. 

- Haftpflicht des Versicherungsnehmers gegenüber dem Geschädigten. 
- Die Eintrittspflicht des Versicherers gegenüber dem Versicherungsnehmer (Deckungspflicht)

Beide Aspekte sind im Streitfall grundsätzlich in verschiedenen Prozessen zu klären. 
Damit soll gesichert werden, dass der vertraglich vereinbarte Versicherungsschutz ungekürzt erbracht wird. 

Dieses Prinzip erfährt aber eine gewichtige Durchbrechung: 
Die im Haftpflichtprozess zwischen dem Geschädigten und dem Versicherungsnehmer getroffenen Feststellungen zur Schadensersatzhaftung entfalten 
Bindungswirkung für eine eventuelle nachfolgende Klärung der Versicherungspflicht im Streitverfahren des Deckungsprozess.
Dies soll verhindern das Fragen die im Haftpflichtprozess geklärt wurden nochmals im Deckungsprozess aufgeworfen werden.

Beispiele für die Bindungswirkung sind die
- materielle Bindung
- Voraussetzungsidentität, d.h. es müssen nur Tatbestände für Schadenersatz erfüllt werden (e.g. §823 BGB) alle nicht 
- direkt damit zusammenhängenden (e.g. Strafbestand) sind unerheblich.
- Haftungstatbestand (eine Haftung ist ein Schadensersatz und wenn dieser in eine Erfüllung umgewandelt wird, verfällt der Deckungsprozess)
- Anspruchsbegründung im Haftungsverhältnis, die Haftung kann nicht bejaht werden nur weil der Schädiger freiwillig oder durch gesetzlichen Zwang
  versichert ist, er muss diese vorher abgeschlossen haben.


Gibt es einen Direktanspruch des Geschädigten gegen den Versicherer?

Nein! Der Haftpflichtversicherungsvertrag ist kein Vertrag zugunsten Dritter i.S.v. § 328 I BGB.

Zur Erleichterung der Anspruchsdurchsetzung lässt § 115 I VVG jedoch Ausnahmen zu, die man als gesetzlich Schuldbeitritt begreifen kann. 
Das betrifft zum einen die Pflichtversicherung nach § 1 PflVG, also die Kfz-Haftpflicht. 
Für den Direktanspruch enthält § 3a PflVG ergänzende Regelungen. Er hat immense praktische Bedeutung. 
Im Verkehrsunfallprozess wird nahezu immer der Haftpflichtversicherer verklagt, daneben oft auch Fahrer als Gesamtschuldner (§ 115 I 4 VVG).

Darüber hinaus besteht ein Direktanspruch insbesondere bei Insolvenz des Versicherungsnehmers bzw. 
Ablehnung des Insolvenzverfahrens mangels Masse (§ 115 I 1 Nr. 2 VVG). 
Dies erfolgt mit Rücksicht auf Verbraucherschutzgesichtspunkten.

Die direkte Inanspruchnahme des Versicherers stellt einen deliktsrechtlichen Schadensersatzanspruch sui generis dar, 
der auf Geldzahlung gerichtet ist (§ 115 I 3 VVG). Sein Umfang richtet sich nach dem Versicherungsverhältnis, wobei
den Geschädigten weder ein Selbstbehalt noch ein Prämienrückstand (§ 114,121).

Selten kann der Geschädigte eine Feststellungsklage gegen VR ausüben, 
zur Prüfung ob eine Freistellung des Schädigers bedingungsgemäß erfolgt ist. 
Dies gilt insbesondere bei Untätigkeit des VN.


Definieren Sie Freistellungs-und Abwehransprüche

Der Freistellungsanspruch ist auf Befreiung von der auf dem Versicherungsnehmer lastenden Haftpflichtschuld, 
mithin auf Vermögensentlastung gerichtet. 
Er ist einzig gegenüber dem Versicherungsnehmer zu erfüllen, während ihm Haftpflichtversicherten eine Leistungsbewirkung i.S.v. § 267 I BGB erfolgt. 
Auf diese Weise ist jeder geschädigte Dritte vor dem Zugriff anderer Gläubiger des Versicherungsnehmers geschützt (vgl. auch § 108 I VVG). 
Als berechtigt gelten die Ansprüche des Dritten dann, wenn der Versicherungsnehmer aufgrund Gesetzes, rechtskräftigen Urteils, (Ziff. 5.1 AHB). 
Der Freistellungsanspruch ist binnen zwei Wochen nach verbindlicher Feststellung der Haftpflicht zu erfüllen (§ 106 S. I VVG, Ziff. 5.1 AHB).

Ihr Umfang nach ist die an den Dritten zu erbringende Entschädigungszahlung auf die vertraglich vereinbarte Versicherungssumme beschränkt 
(Ziff. 6.1 AHB). 

Er ist im Übrigen nicht gehindert, den Freistellungsanspruch an den geschädigten Dritten abzutreten (§ 108 I B VVG, Ziff. 8.1 AHB). 
Das Trennungsprinzip soll insoweit nicht stehen und steht der Abtretung auch nicht entgegen. 
Durch sie wandelt sich der Anspruch in der Person des Dritten zu einem Zahlungsanspruch (im „Direktanspruch durch den Hinterm“). 
Wird dieser (aus abgetretenem Recht) klageweise geltend gemacht, verschwimmen Deckungs- und Haftungsprozess mit zum einemander. 
Dem Versicherer bleiben die Einwendungen aus dem Deckungsverhältnis erhalten (§ 404 BGB).


Definieren Sie den Regress und Zession im juristischen Sinne

Zunächst ist der Begriff des Aufwendungsersatzes zu klären.
Beim Aufwendungsersatz gibt es aber zwei Personen: Denjenigen, der Aufwendungen gemacht hat, und den, 
der daraus einen Vorteil zieht und deshalb die Aufwendungen zu erstatten hat (§§ 284-292 BGB).

E.g. Käufer (K.) bestellt ein Auto beim Verkäufer (V.) und kauft sich ein Navigationssystem kurz vor Lieferung (angegeben von V).
V. liefert jedoch nicht und auch nicht nach einer Fristsetzung.

Der Regress, auch Rückgriff genannt, ist ein Spezialfall des Aufwendungsersatzes. 
Hier sind 3 Parteien gegeben (wie in jeder Kreditsicherungen)
- Der Gläubiger, also derjenige, der den Kredit zur Verfügung stellt.
- Der Schuldner, also derjenige, der den Kredit erhält und
- Der Sicherungsgeber, also der, der mit einem Instrument des Kreditsicherungsrechts die Forderung absichert.

Z.b. kann ein Sicherungsgeber ein Arbeigeber sein, die Schuldner seine Arbeitnehmer und der Gläubiger eine Bank.

Im Falle des Regress geht typischerweise die Forderung vom Gläubiger auf den Sichersgeber über, sprich 
er geht in Regress mit dem Schuldner.

Im Spezialfall von Hafptlichtversicherungen wird der Übergang von der Forderung des Geschädigten an den Schädiger,
auf VR zu Schädiger auch Zession genannt.


Wer will was, von wem, woraus im medizinischen Produkthaftungsgesetz?

- Vertragliche Ansprüche
- Quasivertragliche Ansprüche (Verletzung vorvertraglicher Pflichten/Culpa contrahendo)
- Delikt (z.B. § 823 Abs. 1 BGB, § 1 ProdHaftG)

Nach Art. 10 Abs. 16 Verordnung (EU) 2017/745 sind Medizinprodukther. verpflichtet
Haftpflichtversicherung im ausreichenden Maße abzuschließen.

Insbesonere hat ein Hersteller Schadensersatzpflicht, das kann auch einen Arzt beinhalten 
der bei einer OP die Bestandteile unterschiedlicher Prothesenteile zusammenführt, jedoch
nicht wenn er eine bestehende Prothese durch ein identitisches Modell austauscht.


Was sind die Gesichtspunkte der Risikobewertung in Medizinprodukthaftpftver.

- Entwicklungs-/Produktionsbedingungen (z.B. Erfahrung, Qualität der Marktbeobachtung des Unternehmens)
- Marktbeobachtung/Schadenerfahrung (auch tendenziell steigende Anzahl von OPs in Deutschland)
- OP Praxis / Expertenunterstützung (Erkennbarkeit eines Produkt-mangels im
- „verbauten“ Zustand (Stentbruch, Netze), Identifikation, Risiken einer Revisions-OP)
- Attraktivität der Verwendung (z.B. Adipositas und Bewegungsmangel → mehr Schädigungen des Stützapparates → mehr Ersatz)
Art der hergestellten Medizinprodukte; gibt es Langzeitstudien?
- Risikoklasse (haben Tendenz für Kumulrisiko)
- Stückzahlen
- Verbreitung (regional)
- Klagerisiko / Sammelklagen etabliert?
- Einsatzgebiet pro Land

Kritische Produkte sind somit z.B.
- implantierte Defibrillatoren
- Herz-Stents
- Bandscheibenersatz 
- Empfängnisverhütung
generell US-Market Exposure.
Darum werden keine Versicherungen mit globalem Schutz, sondern länderspezifisch abgeschlossen.

Nach Art. 10 Abs. 16 Verordnung (EU) 2017/745


Was sind die unterschiedlichen Definitionen des Versicherungsfalles im Haftpflichtversichersrecht?

Das Gesetz verzichtet in § 100 VVG auf eine klare Definition des Versicherungsfalles. 
Es spricht nur von der Geltendmachung von Ansprüchen eines Dritten „für eine während der Versicherungszeit eintretende Tatsache“. 
Bei solchen Ansprüchen handelt es sich in der Regel um privatrechtliche Schadensersatzansprüche. 
Allgemein sind davon Personen- und Sachschäden sowie sich daraus ergebende Vermögenseinbußen erfasst (Ziff. I.1 AHB). 
Ansprüche auf Erfüllung von Verträgen oder deren Surrogate scheiden aus (Ziff. I.2 AHB), nicht hingegen 
die in Anspruchskonkurrenz stehenden vertraglichen Schadensersatzansprüche, sofern sie nicht lediglich das Erfüllungsinteresse betreffen.

In der Praxis sind folgende Definitionen je nach Fall gebräuchlich.
Kausalprinzip:
Der Versicherungsfall tritt mit der ersten schadensstiftenden Pflichtverletzung 
(z.B. Beratungs- oder Planungsfehler) des VN ein; Anwendungsbeispiel: Berufshaftpflichtversicherung (§ 1 I AVB Vermögensschaden).

Schadensereignisprinzip: 
Der Versicherungsfall knüpft an den Vorgang an, der unmittelbar in eine Schädigung des Dritten einmündet; Anwendungsbeispiel: 
Allgemeine Haftpflichtversicherung (Ziff. 1.1 AHB).

Feststellungsprinzip: 
Hier wird auf der erste nachprüfbare Feststellung (Manifestation) des Versicherungsfalls abgestellt; 
Anwendungsbeispiel: Umwelthaftpflichtversicherung (§ 8 AVB Umwelthaftpflicht).

Ein weiteres Beispiel ist Ziffer 8.3 ProdHB, die Zusammenfassung mehrerer Versicherungsfälle aus Lieferungen von Erzeugnissen mit gleichen
Mängeln aus der gleichen Ursache für Deckungssumme und Selbstbehalte.
